\section{Resultados}
% Se muestran e interpretan los resultados como se obtuvieron

Como podemos observar en la figura~\ref{fig: p100_60} nuestro algoritmo logra encontrar el minimo global 
de la función de Rosenbrock en un número relativamente bajo de generaciones (alrededor de 60 generaciones), 
lo cual es un buen resultado considerando la complejidad de la función y el tamaño de la población (100 individuos).
Además, la convergencia rápida indica que los operadores de selección, cruza y mutación están funcionando de manera 
efectiva para guiar a la población hacia el óptimo global.

Las gráficas nos permiten ver como la población se va acercando al mínimo global (1,1) a medida que avanzan 
las generaciones ademas de contar con el zoom que nos permite observar con mayor detalle el comportamiento
cerca del mínimo el cual es más notorio en las ultmimas generaciones.

Por otro lado en la figura~\ref{fig: p100_max_10} podemos observar que el algoritmo también es capaz de encontrar
el máximo global de la función de Rosenbrock, aunque en este caso la convergencia es un poco más rápida, 
alcanzando el máximo en un par de generaciones (alrededor de 10 generaciones). Podemos atribuir esto a el 
acotamiento de nuestra función en el dominio definido, lo que facilita la búsqueda del máximo.